%----------------------------------------------------------------------------------------
%    PACKAGES AND THEMES
%----------------------------------------------------------------------------------------
\documentclass[aspectratio=169,xcolor=dvipsnames]{beamer}
\usetheme{CleanEasy}
\usepackage[utf8]{inputenc}
\usepackage{lmodern}
\usepackage[T1]{fontenc}
% \usepackage[brazil]{babel}
\usepackage{fix-cm}
\usepackage{amsmath}
\usepackage{mathtools}
\usepackage{listings}
\usepackage{xcolor}
\usepackage{hyperref}
\usepackage{graphicx} % Allows including images
\usepackage{booktabs} % Allows the use of \toprule, \midrule and \bottomrule in tables
\usepackage{tikz}
\usetikzlibrary{positioning, shapes, arrows, calc, decorations.pathreplacing, arrows.meta, backgrounds, patterns, overlay-beamer-styles}
\usepackage{etoolbox}
\usepackage{animate}
\usepackage{textpos}
\usepackage{transparent}

%----------------------------------------------------------------------------------------
%    LAYOUT CONFIGURATION
%----------------------------------------------------------------------------------------
\pdfstringdefDisableCommands{%
  \def\\{}%
  \def\texttt#1{<#1>}%
}

\newif\ifspecialpage
\specialpagefalse
\addtobeamertemplate{title page}{\specialpagetrue}{\specialpagefalse}

\setbeamertemplate{footline}{
    \leavevmode%
    \hbox{%
        \begin{beamercolorbox}[wd=.2\paperwidth,ht=2.25ex,dp=1ex,left]{title in head/foot}%
            \usebeamerfont{title in head/foot}\hspace*{3ex}\textcolor{fg!30}{Solar Physics + AI}
        \end{beamercolorbox}%
        \begin{beamercolorbox}[wd=.65\paperwidth,ht=1.25ex,dp=1ex,center]{title in head/foot}%
            \usebeamerfont{title in head/foot}\hspace*{-5ex}\textcolor{fg!30}{José P. Marchezi}
        \end{beamercolorbox}%
        \begin{beamercolorbox}[wd=.15\paperwidth,ht=1.25ex,dp=1ex,right]{date in head/foot}%
            \usebeamerfont{date in head/foot}\hspace*{-8ex}\textcolor{fg!30}{\today \hspace*{3ex}
            \insertframenumber{} / \inserttotalframenumber\hspace*{2ex}}
        \end{beamercolorbox}}%
    \vskip0pt\hfill%
}

\AtBeginSection[]
{
  \begin{frame}[plain]
    \begin{tikzpicture}[remember picture,overlay]
      % Fundo limpo
      \fill[MediumBlue!5] (current page.north west) rectangle (current page.south east);
      
      % Título da seção
      \node[
        anchor=west, 
        text width=0.6\paperwidth, 
        align=left, 
        font=\Large\bfseries
      ] at ([xshift=-2cm]current page.center) {
        \insertsectionhead
      };
      
      % Linha horizontal
      \draw[thick, DarkBlue!40] 
        ([xshift=-2cm, yshift=-1cm]current page.center) -- ++(\paperwidth/2,0);
    \end{tikzpicture}
  \end{frame}
}
%----------------------------------------------------------------------------------------
%    TITLE PAGE
%----------------------------------------------------------------------------------------
\title{Research Summary 2025}

\author[J. P. Marchezi]{\small{José P. Marchezi\inst{1}}}

\institute[INPE]{\inst{1}National Institute for Space Research, Brasil; \\ \texttt{\scriptsize{\textcolor{fg!50}{jose.marchezi@inpe.br}}}}

\vspace{-2cm}\date{\today}
% Define positions for logos on title page
\titlegraphic{
  \begin{tikzpicture}[remember picture, overlay]
    %EAGE 2025 logo
    % \node[anchor=north west, xshift=-1cm, yshift=-0.3cm] at (current page.north west) {
    %   % \includegraphics[height=1cm]{logos/logoGeofisica.jpg}
    % };

    % EAGE logo
    \node[anchor=north east, xshift=-6.2cm, yshift=-0.4cm] at (current page.north east) {
      \includegraphics[height=0.5cm]{logos/CNPq_v2017_rgb.jpg}
    };
    
    % CEPETRO logo
    \node[anchor=north east, xshift=-0.8cm, yshift=-0.3cm] at (current page.north east) {
      \includegraphics[height=1.1cm]{logos/inpe_logo.png}
    };


    % % IG logo
    % \node[anchor=north east, xshift=-2.5cm, yshift=-0.3cm] at (current page.north east) {
    %   \includegraphics[height=1.1cm]{logos/LOGO_IG_NAME_D.png}
    % };
    
    %UNICAMP logo
    % \node[anchor=north west, xshift=0cm, yshift=-0.3cm] at (current page.north west) {
    %   \includegraphics[height=1cm]{logos/baner_esqww.png}
    % };
  \end{tikzpicture}
}


%----------------------------------------------------------------------------------------


\begin{document}

\begin{frame}[plain]
  \titlepage
\end{frame}

% \begin{frame}[plain]{Contents}
%   \tableofcontents
% \end{frame}




\begin{frame}{Main Research Project}
\begin{columns}
    \column{0.6\textwidth}
    \begin{itemize}
        \item \textbf{The Influence of Particle Precipitation in the South Atlantic Magnetic Anomaly region on the Geomagnetically Induced Currents}
        \begin{itemize}
            \item Rapid fluctuations in of the geomagnetic field Brazil through high-resolution (temporal) magnetometer data. 
            \item Influence of particle precipitation in the South Atlantic Magnetic Anomaly (SAMA) region on Geomagnetic perturbations. 
            \item Estimate the contribution of the different frequencies of geomagnetic variations on the Geomagnetically Induced Currents (GIC).
            \item Develop predictive models to mitigate GIC impacts on low-latitude.
        \end{itemize}
    \end{itemize}
    
    \column{0.3\textwidth}
    \begin{figure}
        \centering
        \includegraphics[width=1.2\textwidth]{Figures/Picture1.png}
    \end{figure}

\end{columns}

\end{frame}


\section{For High Latitudes}


\begin{frame}{MLT and Latitudinal Distribution}
    \begin{columns}
        \column{0.5\textwidth}
            \begin{itemize}
                \item ICME - Initial Phase: $dB_H/dt$ spikes are broadly dispersed across all MLT (SSC/SI).
                \item ICME - Main Phase: Activity becomes highly concentrated in an intense post-midnight/dawn hotspot.
                \item HSS (All Phases): Spikes are consistently dominated by a pre-midnight hotspot (~22 MLT).
                \item ICMEs produce higher intensity spikes (peaking post-midnight) compared to HSS events (peaking pre-midnight).
                \item ICMEs and HSSs couple differently with the magnetosphere.
            \end{itemize}
        \column{0.5\textwidth}
            \begin{figure}
                \centering
                \includegraphics[width=0.7\textwidth]{Figures/combined_grid_B.pdf}
            \end{figure}
            \vspace{-0.7cm}
%\flushright{\tiny{Marchezi, 2025 (\textit{Preprint})}}
    \end{columns}
\end{frame}

\begin{frame}{Solar Cycle Dependency}
    \begin{columns}
        \column{0.5\textwidth}
            \begin{itemize}
                \item Solar Minima: $dB/dt$ events are most prevalent during the storm recovery phase.
                \item Solar Maxima: Activity is broader and more intense across all phases, peaking in the main phase  due to ICME-driven sudden.
                \item Forecasting Implication: Predictive models must account for these cycle-dependent, phase-specific vulnerabilities, such as the heightened $dB/dt$ risk during the recovery phase of storms at solar minimum.
            \end{itemize}
        \column{0.5\textwidth}
            \begin{figure}
                \centering
                \includegraphics[width=0.7\textwidth]{Figures/combined_grid_B_solarCycle.pdf}
            \end{figure}
            \vspace{-0.7cm}
%\flushright{\tiny{Marchezi, 2025 (\textit{Preprint})}}
    \end{columns}
\end{frame}

\begin{frame}{Side Projects}
    \begin{columns}
        \column{0.5\textwidth}
            \begin{itemize}
            \item \textbf{Modeling and Forecasting Radiation Belt Electron Flux: A Hybrid kappa-Maxwellian Distribution and Neural Network Approach}
            \begin{itemize}
                \item Kappa Distribution
                \begin{equation*}
                    f_\kappa(v) = \frac{n}{\left(\pi\kappa\theta^2\right)^{3/2}}\frac{\Gamma(\kappa+1)}{\Gamma(\kappa-1/2)}\left(1+\frac{v^2}{\kappa\theta^2}\right)^{-(\kappa+1)},
                    \end{equation*}
                \item Maxwellian
                    \begin{equation*}\label{eq:maxwellian}
                    f_M(v) = \frac{n}{\left(2\pi v_{th}^2\right)^{3/2}}\exp{\left(-\frac{v^2}{2v_{th}^2}\right)}.
                    \end{equation*}
            \end{itemize}
        \end{itemize}
        \column{0.5\textwidth}
        \begin{figure}
             \centering
             \includegraphics[width=0.8\linewidth]{Figures/kappaDistr.png}
        
         \end{figure}
    \end{columns}
        
 
    
\end{frame}


\begin{frame}{Side Projects}

\begin{figure}
     \centering
     \includegraphics[width=0.5\linewidth]{Figures/fits_fig3_4.jpg}
 \end{figure}

\begin{figure}
     \centering
     \includegraphics[width=0.5\linewidth]{Figures/fits_fig5_5.jpg}

 \end{figure}
\end{frame}


\begin{frame}{Combining the distributions}
\begin{center}

\begin{tabular}{c}
    
   \includegraphics[width = 0.7\textwidth]{Figures/real_data_2014-02-15_B.png}  \\
    \includegraphics[width = 0.7\textwidth]{Figures/fit_data_2014-02-15_B.png}
\end{tabular}
    
\end{center}
%
% \begin{figure}[H]
%     \centering
%     \includegraphics[width = 0.5\textwidth]{Figures/FitAllDay_20140215.jpg}

% \end{figure}

\end{frame}

\end{document}




\begin{frame}
\frametitle{Algorithm for Storm Identification (using SYM-H)}

\begin{itemize}
    \item \textbf{Storm Identification Trigger:} An event is classified as a storm when the SYM-H index reaches \textbf{$-50~\textrm{nT}$} or less.

   
    \item \textbf{Main Phase Start:} Begins when the SYM-H index reaches the quiet levels prior storm. 

      \item \textbf{Initial Phase Start:} Defined as the point where the index reaches quiet levels \textit{before} its local maximum (the peak before the main phase drop).

    
    \item \textbf{Recovery Phase Start:} This phase begins exactly at the \textbf{minimum SYM-H value} (the storm's peak).

    \item \textbf{Recovery Phase End:} The storm is considered over when the index returns to \textbf{20\%} of its minimum value OR reaches a quiet level of \textbf{$-15~\textrm{nT}$}.
\end{itemize}

\end{frame}

\begin{frame}
\frametitle{Storm Classification by Solar Wind Structure}

\begin{itemize}
    \item Storms were separated based on the related solar wind structure.

    \item \textbf{Two Main Classifications:}
    \begin{itemize}
        \item Interplanetary Coronal Mass Ejections (ICMEs)
        \item High-Speed Streams (HSSs)
    \end{itemize}

    \item \textbf{Identification Sources:}
    \begin{itemize}
        \item Public catalogs:
        \begin{itemize}
            \item Cane and Richardson Lists
            \item DONKI Database
        \end{itemize}
        \item Comparison with other published event lists (Walach \& Grocott, 2019; Pedersen et al., 2024; Juusola et al., 2023; Marchezi et al., 2022)
        \item Visual analysis of solar wind data (NASA/GSFC OMNIWeb)
    \end{itemize}

    \item \textbf{Final Event Count:}
    \begin{itemize}
        \item \textbf{Total: 623} events
        \item ICME-driven: 308
        \item HSS-driven: 315
    \end{itemize}
\end{itemize}

\end{frame}

\begin{frame}{Example storms}
    \begin{tikzpicture}[remember picture, overlay]
    \node[anchor=south west,inner sep=0] at (1.5,-4) {\includegraphics[width=0.8\textwidth]{Figures/stormsICMEHSS.png}};
            
            %\node[anchor=south west,inner sep=0] at (11,-4.1) {\tiny{Marchezi, 2025 (\textit{Preprint})}};
  \end{tikzpicture}

\end{frame}


\section{Results}
\begin{frame}{Superposed Epoch Analysis}
    \begin{columns}
        \column{0.5\textwidth}
            \begin{figure}
                \centering
                \includegraphics[width=0.75\textwidth]{Figures/OMNI_SEA_dBdteventsICME.png}
            \end{figure}
        \column{0.5\textwidth}
            \begin{figure}
                \centering
                \includegraphics[width=0.75\textwidth]{Figures/OMNI_SEA_dBdteventsHSS.png}
            \end{figure}
    \end{columns}
    
\vspace{-0.2cm}
%\flushright{\tiny{Marchezi, 2025 (\textit{Preprint})}}
\end{frame}

\begin{frame}{Dataset}
\begin{columns}
    \column{0.4\textwidth}
        \begin{itemize}
        \item The geomagnetic data is from SuperMag.
            \begin{itemize}
                \item Magnetic coordinate system (N), magnetic east (E), and vertical down (Z) 
            \end{itemize}
        \item All available stations from 1995 to 2024.
        \item SYM-H and solar wind data from OMNIWeb.
    \end{itemize}   
    \column{0.6\textwidth}
    \begin{figure}
        \centering
        \includegraphics[width=\textwidth]{Figures/good_locs_Supermag.png}
    \end{figure}
\end{columns}
\end{frame}


\begin{frame}{MLT Distribution}
    \begin{columns}
        \column{0.5\textwidth}
            \begin{itemize}
                \item $dB_H/dt$ spike occurrence versus Magnetic Local Time (MLT) for ICME-driven vs. HSS-driven storms.
                \item Four magnitude thresholds (from $>50$ nT/min to $>300$ nT/min).
                \item Two main peaks in the pre-midnight (~22 MLT) and dawn (~5 MLT) sectors.ICME vs. 
                \item ICMEs produce more frequent and higher magnitude spikes than 
                \item HSS spikes are fewer and more confined to the pre-midnight sector, especially at high magnitudes.
            \end{itemize}
        \column{0.5\textwidth}
            \begin{figure}
                \centering
                \includegraphics[width=0.7\textwidth]{Figures/MLT_distribution.png}
            \end{figure}
            \vspace{-0.7cm}
%\flushright{\tiny{Marchezi, 2025 (\textit{Preprint})}}
    \end{columns}
\end{frame}
